% =====================================================================
%  Цель и задание лабораторной работы.
% =====================================================================

\textbf{Цель:} освоить на практике основные принципы создания систем
синтеза речи.

\vspace{1em}

\textbf{Задание (вариант 1)}: разработать автоматизированную систему
синтеза речи со следующей комбинацией признаков варианта:
\begin{itemize}
  \item \textbf{Поддерживаемые языки} -- английский;
  \item \textbf{Предметная область} -- научные статьи по
  computer science.
\end{itemize}

\vspace{1em}

Система синтеза должна обеспечивать следующие минимальные
возможности:
\begin{itemize}
  \item ввод текста, или копирование через буфер обмена, или поддержку
  указателя мыши в любом другом приложении;
  \item воспроизведение сгенерированного речевого сигнала для
  введённого текста;
  \item настройки (выбор) голоса диктора, темпа, громкости и других
  параметров для чтения.
\end{itemize}

В отчёт включены: описание структуры разработанной системы,
технологий и методик; основные алгоритмы и принципы реализации
компонент системы (блок-схемы алгоритмов); результаты тестирования;
результаты анализа полученных данных и предложения по улучшению
работы системы.

\vspace{1em}

Синтез речи реализован не как отдельное приложение, а как ещё один
сценарий использования уже существующей информационно-поисковой
системы (ЛР1--ЛР4): озвучить можно любой текст, с которым система
работает -- документ коллекции из предметной области варианта (научная
статья по computer science), реферат, перевод, а также произвольный
текст, введённый вручную или скопированный из любого другого
приложения через буфер обмена. Для синтеза выбрана нейросетевая
библиотека с открытым исходным кодом Piper, работающая локально, без
обращения к облачным сервисам (см. раздел «Технологии и методики»).
